% Part-sized gap in the ToC, so the chapter does not read as part of Part III.
\addtocontents{toc}{\protect\addvspace{2.25em}}
\chapter{Conclusions}
\label{chap:conclusions}

% TODO(canovaccio): expand into full prose.
\lipsum[1]

\section{Answers to the Research Questions}
\label{sec:rq-answers}

% One subsection per question, each in the same four moves: (1) the question in
% one sentence; (2) the answer as a claim -- a declarative sentence that could
% be false, not a summary of what the chapter did; (3) the specific result
% backing it, named; (4) the boundary of the claim. "Yes, under hypotheses (i)
% and (ii) of \cref{thm:deployment-independence}" is an answer; "the thesis
% explored how pulverisation enables flexible deployment" is not one.
%
% RQ1 -> Chapter 5: \cref{thm:deployment-independence}, confirmed in simulation
% by \cref{subsec:pulverization-validation}. \Cref{chap:demonstrator} adds
% scoped realisability evidence -- one admissible deployment, built on physical
% robots -- and nothing more than that. It is not end-to-end validation:
% \cref{sec:emerge-discussion} states that it exercises none of
% \crefrange{chap:deployment-and-reconfiguration}{chap:heterogeneous-gnn}, uses
% neither \cref{chap:language-and-coordination} nor
% \cref{chap:language-based-macroprogramming}, and reports no measurements. Do
% not contradict that statement here.
%
% RQ2 -> Chapter 6: CaMiL's soundness result over the three paradigms, its 46
% use cases, and ScalaTropy's communication primitives -- with the price both
% designs pay, refusing programs the paradigms they unify would accept.
% RQ2.1 -> Chapter 7: the abstraction ladder decides the outcome, not model
% scale.
%
% RQ3 -> Chapters 8-9: three hand-written policies, each reducing the quantity
% it targets at a bounded secondary cost and each needing a person to state what
% a good deployment is, then a learned policy that removes that requirement,
% given a representation that keeps the continuum's heterogeneity and a
% collective summary.
\lipsum[1]

\section{Limitations}
\label{sec:limitations}

% Aggregate the limits the chapters already state, grouped by research question
% rather than restated chapter by chapter: three in
% \cref{sec:pulverization-discussion}, three in
% \cref{sec:llm-macroprogramming-discussion}, four in
% \cref{sec:idhgql-discussion}, and the gaps in \cref{sec:emerge-discussion}.
% Grouping them by question shows which answer is weakest -- RQ3, whose evidence
% is entirely simulated.
\lipsum[1]

\section{Future Work}
\label{sec:future-work}

% Four steps that would close the distance between the demonstrator and the rest
% of the thesis are named in \cref{sec:emerge-discussion}, which promises they
% are taken up here: a local runtime instance on every robot with localisation
% decentralised through radio ranging, communication delays injected into the
% simulation engine so the mixed deployments have a cost, the perturbation
% experiments repeated under controlled conditions, and the collective
% behaviours generated from a natural-language specification. That promise has
% to be kept. Also open, from \cref{sec:idhgql-discussion}: multi-objective
% methods that learn a set of policies covering the Pareto front instead of one
% policy per weight vector, distinguishing application devices by profile, and
% placing bound components jointly rather than independently.
\lipsum[1]
